\textbf{COLUNI 2019 -- Questão 04}\par\medskip \begin{center}\begin{verse}
No descomeço era o verbo.\\
Só depois é que veio o delírio do verbo.\\
O delírio do verbo estava no começo, lá onde a\\
criança diz: \textit{Eu escuto a cor dos passarinhos.}\\
A criança não sabe que o verbo escutar não funciona\\
para cor, mas para som.\\
Então se a\\
criança muda a função de um verbo, ele\\
delira. E pois.\\
Em poesia que é voz de poeta, que é a voz de fazer\\
nascimentos ―\\
O verbo tem que pegar delírio.
\end{verse}
(BARROS, Manoel de. \textbf{Meu quintal é maior do que o mundo}. Rio de Janeiro: Alfaguara, 2015, p. 83.)\end{center}

Releia esse verso: “No \underline{des}começo era o verbo”.

\par

Marque a alternativa em que o prefixo grifado apresenta sentido diferente do observado no excerto:\begin{enumerate}
\item “Meu \underline{des}nome é Andaleço.” (BARROS, Manoel de. \textbf{Meu quintal é maior do que o mundo}. Rio de Janeiro: Alfaguara, 2015, p.103.)
\item “O que não sei fazer \underline{des}conto nas palavras.” (BARROS, Manoel de. \textbf{Meu quintal é maior do que o mundo}. Rio de Janeiro: Alfaguara,
       2015, p. 72.)
\item “As palavras continuam com os seus \underline{des}limites.” (BARROS, Manoel de. \textbf{Meu quintal é maior do que o mundo}. Rio de Janeiro: Alfaguara,
       2015, p. 112.)
\item “Há um \underline{des}agero em mim de aceitar essas medidas.” (BARROS, Manoel de. \textbf{Memórias inventadas}. RJ: Alfaguara, 2018, p. 43)\end{enumerate}
